% !PRINTED 232
\weberpage{232}

%% [running head, not body text -- commented out 2026-08-07] Page furniture (running
%% head, folio, signature) is excluded from this edition by standing policy.
%% 232 \hfill Sixth Section. \hfill \S.\,74.

%<<CONT
we assume that another of the three quantities $B_{0,3}, B_{1,3}, B_{2,3}$
is different from zero.

\indent In order now to treat the equations (5), (6) in a symmetric
way, and in particular to recognise which square root is required
for their solution, we proceed as follows. For brevity we set
\weq{7}{
\begin{array}{rcl}
B_0 &=& \alpha_0\,B_{0,0} + \alpha_1\,B_{0,1} + \alpha_2\,B_{0,2} \\
B_1 &=& \alpha_0\,B_{1,0} + \alpha_1\,B_{1,1} + \alpha_2\,B_{1,2} \\
B_2 &=& \alpha_0\,B_{2,0} + \alpha_1\,B_{2,1} + \alpha_2\,B_{2,2}\,.
\end{array}
}
Then we can represent the equations (5) and (6) as follows:
\weq{8}{
\begin{array}{rcccccl}
\alpha_0\,B_0 & + & \alpha_1\,B_1 & + & \alpha_2\,B_2 & = & 0 \\
\alpha_0\,B_{0,3} & + & \alpha_1\,B_{1,3} & + & \alpha_2\,B_{2,3} & = & 0\,.
\end{array}
}
We multiply the second by an indeterminate factor
$\alpha_3$ and add it to the first, whereby we obtain
\[
\text{(9)} \qquad \alpha_0\,(B_0 + \alpha_3\,B_{0,3}) + \alpha_1\,(B_1 + \alpha_3\,B_{1,3}) + \alpha_2\,(B_2 + \alpha_3\,B_{2,3}) = 0\,,
\]
an equation which is satisfied if we set
\weq{10}{
\begin{array}{rcl}
B_0 + \alpha_3\,B_{0,3} &=& 0 \\
B_1 + \alpha_3\,B_{1,3} &=& \sigma\,\alpha_2 \\
B_2 + \alpha_3\,B_{2,3} &=& -\,\sigma\,\alpha_1\,.
\end{array}
}
From these three equations, in combination with the equation
(6), we have to determine the unknowns $\alpha_0, \alpha_1, \alpha_2, \alpha_3, \sigma$.
We first write the equations out here in full
\weq{11}{
\begin{array}{rcrcrcrcl}
\alpha_0\,B_{0,0} & + & \alpha_1\,B_{0,1} & + & \alpha_2\,B_{0,2} & + & \alpha_3\,B_{0,3} & = & 0 \\
\alpha_0\,B_{1,0} & + & \alpha_1\,B_{1,1} & + & \alpha_2\,(B_{1,2} - \sigma) & + & \alpha_3\,B_{1,3} & = & 0 \\
\alpha_0\,B_{2,0} & + & \alpha_1\,(B_{2,1} + \sigma) & + & \alpha_2\,B_{2,2} & + & \alpha_3\,B_{2,3} & = & 0 \\
\alpha_0\,B_{3,0} & + & \alpha_1\,B_{3,1} & + & \alpha_2\,B_{3,2} & & & = & 0\,,
\end{array}
}
and if we eliminate $\alpha_0, \alpha_1, \alpha_2, \alpha_3$ from these, there results
a quadratic equation for $\sigma$, after whose solution the
ratios $\alpha_0 : \alpha_1 : \alpha_2 : \alpha_3$ can be determined from linear equations.

\indent The quadratic equation for $\sigma$, however, reads in determinant form
\[
\begin{vmatrix}
B_{0,0}, & B_{0,1}, & B_{0,2}, & B_{0,3} \\
B_{1,0}, & B_{1,1}, & B_{1,2} - \sigma, & B_{1,3} \\
B_{2,0}, & B_{2,1} + \sigma, & B_{2,2}, & B_{2,3} \\
B_{3,0}, & B_{3,1}, & B_{3,2}, & 0
\end{vmatrix} = 0\,.
\]
